% =========================================================================
\section{The Structural Score Across Scales}
\label{app:instrument}

Table~\ref{tab:scale} repeats the evaluation of Table~\ref{tab:instrument} on
the larger LINK graphs. The full cleaned graph is not an order of magnitude
larger than the $5\times10^{4}$ subgraph; the scale effect reported in the
first version of this paper, where the largest graph had \V{v1.link.nodes.64w}
nodes, came mostly from the hundreds of thousands of single-receipt airdrop
recipients that the cleaning rule removes.

\begin{table}[H]
\centering\small
\caption{AUC for the core label on LINK at three scales. GraphSAGE: mean
$\pm$ s.d.\ over 10, 5 and 3 initialisations respectively.}
\label{tab:scale}
\setlength{\tabcolsep}{3.5pt}
\begin{tabular}{@{}lccc@{}}
\toprule
Scale (nodes) & Centrality & Structural & GraphSAGE \\
\midrule
$10^{4}$ (\V{link.nodes.1w}) & \V{c1.link.cent} & \V{c1.link.struct} & \V{c1.link.gnn} \\
$5\times10^{4}$ (\V{link.nodes.5w}) & \V{c1.link.5w.cent} & \V{c1.link.5w.struct} & \V{c1.link.5w.gnn} \\
Full (\V{link.nodes.full}) & \V{c1.link.full.cent} & \V{c1.link.full.struct} & \V{c1.link.full.gnn} \\
\bottomrule
\end{tabular}
\end{table}

\paragraph{Neighbourhood purity}
For each core node we compute the fraction of its ten nearest neighbours in
the representation that are also core. Table~\ref{tab:purity} reports the
absolute purity rather than its ratio to the random baseline, because the
baseline (the positive rate) varies with graph size. The GraphSAGE embedding
places core nodes near one another more consistently than the structural
features do.

\begin{table}[H]
\centering\small
\caption{Absolute neighbourhood purity of core nodes ($k=10$).}
\label{tab:purity}
\setlength{\tabcolsep}{3.5pt}
\begin{tabular}{@{}lcccc@{}}
\toprule
 & Centrality & Structural & GraphSAGE & Ratio \\
\midrule
LINK $10^{4}$ & \V{c1.link.1w.pur.cent} & \V{c1.link.1w.pur.struct} & \V{c1.link.1w.pur.gnn} & \V{c1.link.1w.pur.ratio} \\
LINK $5\times10^{4}$ & \V{c1.link.5w.pur.cent} & \V{c1.link.5w.pur.struct} & \V{c1.link.5w.pur.gnn} & \V{c1.link.5w.pur.ratio} \\
LINK full & \V{c1.link.full.pur.cent} & \V{c1.link.full.pur.struct} & \V{c1.link.full.pur.gnn} & \V{c1.link.full.pur.ratio} \\
UNI $10^{4}$ & \V{c1.uni.1w.pur.cent} & \V{c1.uni.1w.pur.struct} & \V{c1.uni.1w.pur.gnn} & \V{c1.uni.1w.pur.ratio} \\
\bottomrule
\end{tabular}
\end{table}

\paragraph{Predicting large balances}
With balance removed from the inputs, the three representations predict
membership in the top 5\% by balance at similar levels on the $10^{4}$
subgraphs (LINK: \V{c1.link.1w.bal.cent}, \V{c1.link.1w.bal.struct},
\V{c1.link.1w.bal.gnn} for centrality, structural features and GraphSAGE;
UNI: \V{c1.uni.1w.bal.cent}, \V{c1.uni.1w.bal.struct},
\V{c1.uni.1w.bal.gnn}). In the first version of this paper the GraphSAGE
embedding was clearly better at this task on LINK; that advantage depended
on the airdrop transfers (\ref{app:cleaning}).

\paragraph{Roles}
Clustering the GraphSAGE embedding into eight groups ($k$-means) yields on
both ledgers one small cluster that is highest in both core rate and median
balance---\V{link.core.size} LINK nodes with a core rate of \V{link.core.rate}
and a median of \V{link.core.medcp} counterparties, \V{uni.core.size} UNI
nodes with \V{uni.core.rate} and \V{uni.core.medcp}---above a periphery of
low-degree addresses (Table~\ref{tab:typology}). The infrastructure cluster
is \V{link.core_share} of the LINK subgraph and \V{uni.core_share} of the UNI
subgraph.

\begin{table*}[t]
\centering\footnotesize
\caption{Clusters of the GraphSAGE embedding on the $10^{4}$ subgraphs, sorted
by core rate. Cluster indices are arbitrary and the two ledgers are aligned by
rank only.}
\label{tab:typology}
\setlength{\tabcolsep}{4.5pt}
\begin{tabular}{@{}rrrr c rrrr@{}}
\toprule
\multicolumn{4}{c}{LINK} & & \multicolumn{4}{c}{UNI} \\
\cmidrule(r){1-4}\cmidrule(l){6-9}
Size & Core rate & Med.\ cp. & Med.\ bal. & & Size & Core rate & Med.\ cp. & Med.\ bal. \\
\midrule
\V{tab.typology}
\bottomrule
\end{tabular}
\end{table*}

% =========================================================================
\section{Cleaning Specifications}
\label{app:cleaning}

\paragraph{Sensitivity of the airdrop rule}
With the amount threshold at 0.01 LINK, $N = 100$, $500$ and $1{,}000$ flag
\V{link.clean.n100.senders}, \V{link.clean.n500.senders} and
\V{link.clean.n1000.senders} senders and remove \V{link.clean.n100.tx},
\V{link.clean.n500.tx} and \V{link.clean.n1000.tx} transfers. The additional
senders at $N=100$ are small dust broadcasters spread across the window.
Raising the amount threshold to 0.1 at $N = 1{,}000$ changes nothing
(\V{link.clean.amt0.1.senders} senders, \V{link.clean.amt0.1.tx} transfers).
On UNI the rule flags \V{uni.clean.n1000.senders} senders at $N = 1{,}000$
and $N = 500$, and \V{uni.clean.n100.senders} at $N=100$.

\paragraph{Data specifications}
Table~\ref{tab:specdata} reruns the $10^{4}$ analyses under each cleaning
specification, and with repeated pairs overwritten instead of summed. The
airdrop transfers are what made GraphSAGE appear better than the structural
features at predicting large LINK balances in the first version; with them
removed the two are level. On UNI, where the airdrop rule removes nothing,
the ``airdrop kept'' and main specifications coincide, as do ``zero-value
kept'' and ``both kept''; each pair was computed in a separate run and the
results are identical, which also checks determinism.

\begin{table*}[t]
\centering\footnotesize
\caption{The $10^{4}$ analyses under alternative data specifications. AUC:
core label (structural features, GraphSAGE) and top-5\% balance
(GraphSAGE). Brokers: size of the hidden-broker set / overlap with the main
set.}
\label{tab:specdata}
\setlength{\tabcolsep}{3.5pt}
\begin{tabular}{@{}l rrrrr rrrrr@{}}
\toprule
& \multicolumn{5}{c}{LINK} & \multicolumn{5}{c}{UNI} \\
\cmidrule(lr){2-6}\cmidrule(l){7-11}
Specification & Nodes & Struct. & GNN & Bal.\ GNN & Brokers & Nodes & Struct. & GNN & Bal.\ GNN & Brokers \\
\midrule
\V{tab.specdata}
\bottomrule
\end{tabular}
\end{table*}

\paragraph{All combinations}
Table~\ref{tab:combos} crosses the five holder populations, three address
types and two units of analysis of Table~\ref{tab:spec}.

\begin{table*}[t]
\centering\footnotesize
\caption{Balance concentration over the full grid of specifications.
``Active'': transfer-active after cleaning; ``Active, airdrop kept'': with the
dust-airdrop transfers counted as activity (zero-value transfers removed);
$\geq$ 0.01 and $\geq$ USD 100: balance floors in tokens and in value at the
snapshot price.}
\label{tab:combos}
\setlength{\tabcolsep}{4pt}
\begin{tabular}{@{}lll rr rr@{}}
\toprule
& & & \multicolumn{2}{c}{LINK} & \multicolumn{2}{c}{UNI} \\
\cmidrule(lr){4-5}\cmidrule(l){6-7}
Population & Address type & Unit & HHI & Gini & HHI & Gini \\
\midrule
\V{tab.combos}
\bottomrule
\end{tabular}
\end{table*}

% =========================================================================
\section{The Matched Control}
\label{app:control}
Centralized-exchange addresses: one blind item, a Gate.io deposit address,
was recorded with its exchange name, and another, an HTX hot wallet, as
unknown. Before unblinding we fixed the rule that exchange deposit and
hot-wallet addresses are not routing protocols. Scoring the Gate.io item as
named instead changes the unmatched group from \V{ctl.uniform_k} to
\V{ctl.uniform_k_cex} named of ten and leaves every other figure unchanged.
The blind list, the group key and the scoring script are archived; the
commit that fixes the rule precedes the commit that unblinds.

% =========================================================================
\section{Reproducibility}
\label{app:repro}

\paragraph{Balances and data} Balances are rebuilt from transfers to the
snapshot block and checked against \texttt{totalSupply()} and a sample of
\texttt{balanceOf()} calls at that block; the LINK genesis allocation is the
only correction needed (Section~\ref{sec:data}). The transfer records are
identical, record for record, to those used in the first version.

\paragraph{Prices and fee contracts} The USD~100 floor uses
\texttt{latestRoundData()} of the Chainlink LINK/USD feed
\V{link.feed} (round \V{link.feed_round}, answer \V{link.feed_answer} at
eight decimals) and UNI/USD feed \V{uni.feed} (round \V{uni.feed_round},
answer \V{uni.feed_answer}), both read at the snapshot block. The
\texttt{TokenJar}, releaser and threshold of the fee mechanism are read on
Ethereum at the first and last block of the window and on
\V{uni.l2.chains} at the block closest to the end of the window; the
addresses come from an archived copy of the deployment page
\cite{uniswap2026deployments} and are checked against each jar's
\texttt{releaser()}.

\paragraph{Edge weights} Repeated ordered pairs must be summed
(Section~\ref{sec:edgecase}).

\paragraph{Labels} All label-dependent results use the archived snapshot.
Its 29{,}945 entries become 29{,}772 addresses after case normalisation; the
first version of this paper described this difference as drift in the
upstream compilation, which was a mistake on our part.

\paragraph{Token-transfer view} Manual verification must use Etherscan's
ERC-20 token-transfer view: a routing contract called by other contracts
records few transactions of its own. Etherscan name tags also change; one
router selected here was tagged ``Uniswap V4: Universal Router'' when first
verified and ``Uniswap: Universal Router v2'' on 6 October 2026.

\paragraph{Determinism} The first version attributed run-to-run variation in
the embeddings to non-deterministic GPU scatter operations. The main cause
was that PyTorch Geometric's \texttt{negative\_sampling} draws from Python's
\texttt{random} module, which seeding PyTorch does not reach. We now seed
Python, NumPy and PyTorch, train on CPU with deterministic algorithms, give
gensim a stable hash function so that node2vec does not depend on the
interpreter's hash seed, break ties in the activity ranking by address, and
insert graph nodes in address order so that sampled betweenness does not
depend on the order in which records are read. Two independent runs of the
pipeline produce byte-identical results. gensim emits an ignored-exception
message from its BLAS dot-product routine during node2vec training; outputs
are nonetheless deterministic and consistent with the first version.

\paragraph{Numbers} Every number in the text and tables is generated from the
archived result files by one script, which also checks each qualitative
statement in the text against the numbers it rests on.

% =========================================================================
\section{Changes from the First Version}
\label{app:changes}
This version (arXiv v2) corrects and extends the first (v1).

\paragraph{Facts and references}
Four references in v1 did not correspond to the cited works (Gudgeon et al.,
Campajola et al., Zarir et al., Grajales et al.) and have been replaced by
works we checked against their texts; several DOIs and links were corrected.
Uniswap's proposal threshold is $10^{6}$ UNI of delegated votes, not
$2.5\times10^{6}$ as stated in v1; we read it, and the quorum, from the
governor contract. The protocol fee that v1 described as a future possibility
had been activated in December 2025; Section~\ref{sec:governance} now treats
it as the central case. Benchmark HHIs without a checkable source (Canadian
and Swedish banking, LPMCL gold clearing, a card-network figure, a U.S.\
banking figure, and the Gini benchmarks of a Lorenz figure) were removed or
replaced by figures computed from public data.

\paragraph{Data and method}
Balances are now a snapshot at the end of the window for both tokens
(Section~\ref{sec:data}); v1 used extraction-date balances. Zero-value
transfers and a coordinated dust airdrop are removed by stated rules, which
changes the LINK subgraphs and removes v1's apparent scale effect. Training
is now deterministic (\ref{app:repro}). The rule layer is measured in
delegated votes rather than balances. Table~\ref{tab:spec} gains rows for
balance floors, and \ref{app:cleaning} gives the full combination grid.

\paragraph{Findings revised}
The Gini coefficient is insensitive to choices about the top of the
distribution but not to balance floors (Section~\ref{sec:grid}). The
hidden-broker sets are \V{link.brokers} and \V{uni.brokers} addresses rather
than \V{v1.link.brokers} and \V{v1.uni.brokers}, with \V{shared.brokers}
rather than \V{v1.shared.brokers} shared; the label compilation they are
defined against dates from 2023 (Section~\ref{sec:verify}). The difference
between matched and unmatched controls reported as significant in v1 is not
significant here. v1's attribution of a cluster-level reversal to edge
aggregation is withdrawn (Section~\ref{sec:edgecase}); its description of a
change in the label compilation as upstream drift was a case-normalisation
artefact of our own; its explanation of non-reproducible embeddings was
wrong (\ref{app:repro}); and its finding that the learned embedding
predicts large balances better than structural features depended on the
airdrop transfers.

\paragraph{Structure}
The research questions are reduced to three, matched to the three
contributions; the specification table is in the main text with a reporting
checklist; the matched control follows the definition of the criterion.
